\documentclass[11pt,a4paper]{article}

\usepackage{fontspec}
\usepackage[magyar]{babel}
\usepackage{amsmath,amssymb,amsthm}
\usepackage{hyperref}
\usepackage{faktor}

\hypersetup{
  colorlinks=true,
  linkcolor=blue,
  urlcolor=blue,
  citecolor=blue
}

\newtheorem{tetel}{Tétel}[section]
\newtheorem{allitas}[tetel]{Állítás}
\newtheorem{lemma}[tetel]{Lemma}
\newtheorem{kovetkezmeny}[tetel]{Következmény}

\theoremstyle{definition}
\newtheorem{definicio}[tetel]{Definíció}
\newtheorem{pelda}[tetel]{Példa}
\newtheorem{megjegyzes}[tetel]{Megjegyzés}

\renewcommand{\proofname}{Bizonyítás}

\DeclareMathOperator{\Frac}{Frac}
\DeclareMathOperator{\Spec}{Spec}
\DeclareMathOperator{\ord}{ord}
\DeclareMathOperator{\length}{length}

\newcommand{\OO}{\mathcal O}
\newcommand{\mm}{\mathfrak m}
\newcommand{\pp}{\mathfrak p}
\newcommand{\qq}{\mathfrak q}
\newcommand{\ZZ}{\mathbb Z}
\newcommand{\NN}{\mathbb N}
\newcommand{\kk}{\Bbbk}

\title{Értékelések és értékelésgyűrűk\\
\large Különös tekintettel a diszkrét értékelésgyűrűkre}
\author{Előadásjegyzet az Algebrai Görbék c. tárgyhoz}
\date{}

\begin{document}

\section{Algebrai görbék és pontjaik}

\begin{definicio}
    A $k$ test felett $K$ egy \emph{algebrai görbe}, a $k(t)$ racionális törtfüggvénytest egy véges bővítése.
\end{definicio}

\begin{definicio}
   $|K|$, azaz \emph{K pontjai} a $\{ \nu : K \rightarrow \mathbb{Z} \cup \{\infty\}\}$ halmaz, ahol $\nu$ normalizált (diszkrét) értékelés, $\nu|_k$ triviális, de $\nu$ nem triviális.
\end{definicio}

\begin{pelda}
    $x^2+y^2=1$ görbe $\mathbb{R}$ vagy $\mathbb{C}$ felett:
    $$Frac\left (\faktor {k[x,y]}{x^2+y^2=1}\right)$$
\end{pelda}

\subsection{A racionális törtfüggvénytest, pontjai}

Mi lehet $|k(t)|?$

\begin{lemma}
Legyen $i: A \to B$ egy gyűrűhomomorfizmus, és legyen $P \vartriangleleft _p~ B$ egy prímideál. Ekkor $i^{-1}(P) \vartriangleleft _p A$.
\end{lemma}

\begin{proof}
Ideál, mert persze $0 \in i^{-1}(P)$
\\
$x, y \in i^{-1}(P)$ $\Rightarrow$ $i(x+y) = i(x) + i(y) \in P$
\\
$x \in i^{-1}(P)$ és $y \in A$ $\Rightarrow$ $i(xy) = i(x)i(y) \in P$, mivel $P$ ideál $B$-ben.
\\~\\
Valódi ideál, mivel $1 \notin i^{-1}(P)$, hiszen a homomorfizmusok egységelemtartók, és $1 \notin P$, mivel $P$ prímideálként valódi
\\~\\
Prím tulajdonságát is teljesül: Tegyük fel, hogy $x, y \in A$ és $xy \in i^{-1}(P)$. Ekkor $i(xy) = i(x)i(y) \in P$. Mivel $P$ prímideál $B$-ben, ezért $i(x) \in P$ vagy $i(y) \in P$, tehát pont $x \in i^{-1}(P)$ vagy $y \in i^{-1}(P)$.
\end{proof}

$\longrightarrow$ a prímideálok jók


TFH. $A \subseteq \OO \subseteq F$, $A$ gyűrű, $\OO$ DVR, $F=Frac(\OO)$
 $$\longrightarrow \exists! \; \; m \vartriangleleft_m \OO \Rightarrow m \cap A \vartriangleleft_p A$$
Legyen $\nu \in |k(t)|$. Mivel $\nu(x)+\nu(x^{-1})=\nu(1)=0$, így $\nu(x)$ vagy $\nu(x^{-1})$ pozitív.
\begin{enumerate}
    \item Első eset: $x \in \OO_\nu$. Ekkor $k[x] \subseteq \OO_{\nu}$
    \item Második eset: $x^{-1} \in \OO_{\nu}$. Ekkor $k[x^{-1}] \subseteq \OO_\nu$
\end{enumerate}

Mivel $k[x]$ (és $k[x^{-1}]$ is) főideálgyűrű, a $k[x] \cap m_\nu$ prímideál egy $(f)$ alakú ideál lesz, ahol $f \in k[x]$ egy irreducibilis polinom.


\begin{allitas}
 $$|k(t)| = \{ \nu \; |\; \OO_\nu = k[x]_{p_\nu}, \, p_\nu \vartriangleleft k[t] \text{ irreducibilis polinom} \} \cup \{ \nu_\infty \}$$
ahol $k(t) \subseteq k((t^{-1})) \longrightarrow^\nu \; \mathbb{Z} \cup \{ \infty \}$.
\end{allitas}

\begin{pelda}
Ha $k$ algebrailag zárt (pl. $k = \mathbb{C}$), akkor $p_\nu=(x-\alpha ): \; \alpha \in \mathbb{C} $ ($a \in \mathbb{C}$). Ekkor
$$|\mathbb{C}(t)| = \mathbb{C} \cup \{\infty\},$$

ami a Riemann-gömb, avagy a $\mathbb{P}^1(\mathbb{C})$ projektív egyenes.
\end{pelda}

\end{document}
